%% kapa: exam=1 round=32 year=2021 period=1 form=A subject=civil_law subject_ko=민법 seq=23 number=23
%% source: Q-Net 감정평가사 기출문제 (cst00309) artlSeq=5208397 file=extracted/1008/5208397/A형/2021년도 제32회 감정평가사 1차 1교시 A형.pdf
%% conversion: kordoc 4.24.51 -> tools/gen1.py (auto); review=proofread-sonnet; flags=none
\begin{question}{23}
총유에 관한 설명으로 옳지 않은 것은? (다툼이 있으면 판례에 따름)
\begin{choices}
  \item 비법인사단이 총유물에 관한 매매계약을 체결하는 행위는 총유물의 처분행위가 아니다.
  \item 비법인사단이 타인 간의 금전채무를 보증하는 행위는 총유물의 관리·처분행위가 아니다.
  \item 총유물의 보존행위는 특별한 사정이 없는 한 구성원이 단독으로 결정할 수 없다.
  \item 비법인사단의 대표자는 총유재산에 관한 소송에서 단독으로 당사자가 될 수 없다.
  \item 비법인사단인 주택조합이 주체가 되어 신축 완공한 건물로서 일반에게 분양되는 부분은 조합원 전원의 총유에 속한다.
\end{choices}
\end{question}
\answer{1}
